\documentclass{beamer}

\usepackage{floatrow}
\usepackage{/home/user1/code/tex/sty}
\setbeamercolor{frametitle}{fg=black}
\setbeamertemplate{navigation symbols}{}

\usepackage{caption}
\captionsetup[figure]{labelformat=empty}

\newcommand{\figpic}[3]{
\begin{figure}[H]
\begin{center}
\includegraphics[height=#1]{#2}
\caption{#3}
\end{center}
\end{figure}}

\newcommand{\pic}[2]{
\begin{center}
\includegraphics[height=#1]{#2}
\end{center}}

\begin{document}

\begin{frame}
\frametitle{Outline}
\begin{enumerate}
\item Graphs
\item Networks
\begin{itemize}
\item Random network
\item Power-law network
\begin{itemize}
\item Friendship Paradox
\end{itemize}
\end{itemize}
\end{enumerate}
\end{frame}

\begin{frame}
\frametitle{Examples of real networks}
\begin{enumerate}
\item Social Networks: Nodes are people, edges are connections.
\end{enumerate}
\end{frame}

\begin{frame}
\frametitle{Types of networks}
\begin{enumerate}
\item Probability Distribution
\begin{enumerate}
\item Exponential Network
\item Random Network
\item Poisson Network
\end{enumerate}
\item Algorithm
\begin{enumerate}
\item BA Network
\item BA Fitness
\end{enumerate}
\end{enumerate}
\end{frame}

\begin{frame}
\frametitle{Properties of a network}
\begin{enumerate}
\item Number of types of nodes, number of types of edges.
\item For each type, number of nodes and number of edges.
\item Average number of edges: $\langle k \rangle$
\item Components
\begin{itemize}
\item Size of the largest component.
\item Connected: only one component.
\end{itemize}
\item Probability that a random selected node has $k$ edges: $p(k)$
\begin{itemize}
\item $\int_{k_{\textrm{min}}}^{k_{\textrm{max}}} p(k) = 1$
\end{itemize}
\end{enumerate}
\end{frame}

\begin{frame}
\frametitle{Random Network}
\begin{itemize}
\item Construction
\begin{enumerate}
\item Start with $N$ nodes and no edges.
\item For each pair of nodes, connect them with probability $p$.
\end{enumerate}
\end{itemize}
\end{frame}

\begin{frame}
\frametitle{Random Network}
\begin{columns}
\column{0.5\textwidth}
\figpic{7cm}{1.pdf}{}
\column{0.5\textwidth}
\figpic{7cm}{2.pdf}{}
\end{columns}
\end{frame}

\begin{frame}
\frametitle{Random Network}
\begin{center}
Construct several networks with equal $N$ but increasing $p$.
\end{center}
\end{frame}


\begin{frame}
\frametitle{Random Network}
\begin{columns}
\column{0.5\textwidth}
% beispiel-11
\pic{4cm}{5.png}
\pic{3cm}{6.pdf}
$N=500$, $p=0.0001$, $\langle k \rangle = 0.03$
\column{0.5\textwidth}
% beispiel-12 
\pic{4cm}{8.png}
\pic{3cm}{7.pdf}
$N=500$, $p=0.001$, $\langle k \rangle = 0.524$
\end{columns}
\end{frame}

\begin{frame}
\frametitle{Random Network}
\begin{columns}
\column{0.5\textwidth}
% beispiel-13
\pic{4cm}{10.png}
\pic{3cm}{9.pdf}
$N=500$, $p=0.003$, $\langle k \rangle = 1.5$
\column{0.5\textwidth}
% beispiel-14
\pic{4cm}{12.png}
\pic{3cm}{11.pdf}
$N=500$, $p=0.01$, $\langle k \rangle = 5.1$
\end{columns}
\end{frame}

\begin{frame}
\frametitle{Random Network}
\begin{columns}
\column{0.5\textwidth}
\figpic{7cm}{3.pdf}{$N=500$, $p=0.01$, $\langle k \rangle = 5.15$}
\column{0.5\textwidth}
\figpic{4cm}{4.png}{}
\end{columns}
\end{frame}

\begin{frame}
\begin{center}
\textit{Component: connected set of nodes}\\~\\
What is the largest component of a random graph?\\~\\
When does a graph become fully connected?\\~\\
How does it scale with $p$?
\end{center}
\end{frame}

\begin{frame}
\frametitle{Random Network}
\begin{columns}
\column{0.5\textwidth}
\pic{4cm}{datum-1.png}
\pic{3cm}{datum-1.pdf}
$N=500$, $p=0.0001$, $N_G=3$
\column{0.5\textwidth}
\pic{4cm}{datum-2.png}
\pic{3cm}{datum-2.pdf}
$N=500$, $p=0.001$, $N_G=9$
\end{columns}
\end{frame}

\begin{frame}
\frametitle{Random Network}
\begin{columns}
\column{0.5\textwidth}
\pic{4cm}{datum-3.png}
\pic{3cm}{datum-3.pdf}
$N=500$, $p=0.003$, $N_G=272$
\column{0.5\textwidth}
\pic{4cm}{datum-4.png}
\pic{3cm}{datum-4.pdf}
$N=500$, $p=0.01$, $N_G=494$
\end{columns}
\end{frame}

\n{kDurch}{\langle k \rangle}
\begin{frame}
\frametitle{Random Network}
\begin{center}
\pic{4cm}{phaseverschiebung.pdf}
\g{\br{1}{N}=\br{1}{500}=0.002}
\g{p_k = \br{1}{N} \Rightarrow \kDurch = 1}
\textit{Phase transition to a single giant component.}
\end{center}
\end{frame}

\begin{frame}
\g{\sigma_k = \sqrt{\langle k \rangle}}
\g{k = \kDurch \pm \sqrt{\kDurch}}
\begin{center}
Nodes never deviate significantly from $\kDurch$.\\~\\
\textit{``scale''}
\end{center}
\end{frame}

\begin{frame}
\frametitle{Types of Networks}
\begin{itemize}
\item Random Network
\item Power Law Network
\end{itemize}
\end{frame}

\begin{frame}
\begin{center}
\g{p_k \propto k^{-\gamma}}
\textit{Power Law Network}
\end{center}
\end{frame}

\n{kneighbor}{k_{\textrm{Neighbor}}}
\n{kND}{\langle \kneighbor \rangle}

\n{kmin}{k_{\textrm{min}}}
\n{kmax}{k_{\textrm{max}}}
\n{emg}{1-\gamma}
\n{zmg}{2-\gamma}
\n{dmg}{3-\gamma}
\n{kme}{\kmax^{\emg}-\kmin^{\emg}}
\n{kmz}{\kmax^{\zmg}-\kmin^{\zmg}}
\n{kmd}{\kmax^{\dmg}-\kmin^{\dmg}}
\begin{frame}
\g{\int p_k dk = 1}
\g{\int_k Ck^{-\gamma} = 1}
\g{C(\br{k^{1-\gamma}}{1-\gamma})_k=1}
\g{\boxed{C=\br{1-\gamma}{\kme}}}
Normalization of a power law network.
\end{frame}

\n{kD}{\langle k \rangle}
\begin{frame}
\g{\kD &= \int k C k^{-\gamma} = C\int k^\emg = C(\br{k^{\zmg}}{\zmg})_{\kmin}^{\kmax}\\
&= \br{\emg}{\kme} \br{\kmz}{\zmg}\\
&= \boxed{\br{\emg}{\zmg} \br{\kmz}{\kme}}}
\begin{center}
\textit{$p_k$ and $\kD$ are independent of the size a power law network.}
\end{center}
\end{frame}

\begin{frame}
\frametitle{Power Law Network}
\begin{itemize}
\item Randomly select a node
\begin{itemize}
\item How many links do you have? $\kDurch$
\end{itemize}
\begin{center}\textit{or}\end{center}
\item Randomly select a link 
\begin{itemize}
\item Go to one of its nodes
\item How many links does that node have? $\langle \kneighbor \rangle$
\end{itemize}
\end{itemize}
\g{\langle k \rangle \neq \langle \kneighbor \rangle}
\begin{center}\textit{``Friendship Paradox''}\end{center}
\end{frame}

\begin{frame}
\begin{center}
\textit{Use $kp_k$ as a probability distribution.}\\~\\
It is most likely to connect to a node which already has $\kDurch$ links.\\~\\
It is $k$ times more likely to connect to a node with $k$ links than a node with 1 link (assuming equal number of $k$-nodes and 1-nodes).
\end{center}
\g{\int_0^{\infty} kp_k \propto \textrm{Average number of links of neighbor.}}
\end{frame}

\begin{frame}
\g{q_k=Akp_k=AkCk^{-\gamma}= \br{\emg }{\kme} Ak^{\emg} }
Plug into normalization.
\g{\int q_k dk = 1}
Solve for $A$.
\g{\boxed{A=\br{\zmg}{\emg} \br{\kme}{\kmz}}}
\end{frame}

\begin{frame}
\g{\kND &= \int kq_k dk = AC \int k^{\zmg} dk = AC(\br{k^{\dmg}}{\dmg})_{\kmin}^{\kmax}\\
&= AC \br{\kmd}{\dmg} \\
&= \boxed{\br{\zmg}{\dmg} \br{\kmd}{\kmz} }}
The expected degree of a neighbor in a Power Law Network.
\end{frame}

\begin{frame}
\g{\kDurch = \br{\emg}{\zmg} \br{\kmz}{\kme} \quad \kND = \br{\zmg}{\dmg} \br{\kmd}{\kmz} }
\begin{center}
\includegraphics[height=5cm]{13.pdf}\\
\textit{``Friendship Paradox''}
\end{center}
\end{frame}

\end{document}
